\documentclass[a4paper,12pt]{article}

\usepackage{amsmath,amssymb}
\usepackage{geometry}
\geometry{margin=2cm}

\title{Formularium Goniometrie}
\author{}
\date{}

\begin{document}

\maketitle

\section*{4.1 Goniometrische cirkel en radialen}

\[
2\pi \text{ rad} = 360^\circ,\qquad
\pi \text{ rad} = 180^\circ,\qquad
\frac{\pi}{2} = 90^\circ
\]

Booglengte:
\[
s = \alpha r
\]

\section*{4.2 Sinus, cosinus, tangens en cotangens}

Punt op de eenheidscirkel:
\[
P=(\cos \alpha,\ \sin \alpha)
\]

Fundamentele formule:
\[
\sin^2 \alpha + \cos^2 \alpha = 1
\]

Tangens en cotangens:
\[
\tan\alpha = \frac{\sin\alpha}{\cos\alpha},\qquad
\cot\alpha = \frac{\cos\alpha}{\sin\alpha}
\]

Symmetrie:
\[
\sin(-\alpha)=-\sin\alpha,\qquad
\cos(-\alpha)=\cos\alpha
\]
\[
\tan(-\alpha)=-\tan\alpha,\qquad
\cot(-\alpha)=-\cot\alpha
\]

Perioden:
\[
\tan(\alpha+\pi)=\tan\alpha,\qquad
\cot(\alpha+\pi)=\cot\alpha
\]

\subsection*{Bijzondere waarden}

\[
\begin{array}{c|ccc}
\alpha & 0 & \frac{\pi}{2} & \pi \\ \hline
\sin\alpha & 0 & 1 & 0 \\
\cos\alpha & 1 & 0 & -1 \\
\tan\alpha & 0 & \text{g.d.} & 0
\end{array}
\]

\subsection*{Goniometrische vergelijkingen}

\[
\sin x=A \quad\Rightarrow\quad
x = B + 2k\pi \ \text{of}\ x=\pi - B + 2k\pi
\]

\[
\cos x=A \quad\Rightarrow\quad
x = B + 2k\pi \ \text{of}\ x=-B + 2k\pi
\]

\[
\tan x=A \quad\Rightarrow\quad
x = B + k\pi
\]

\[
\cot x=A \quad\Rightarrow\quad
x = B + k\pi
\]

\section*{4.3 Verwante hoeken}

Complementair: $\alpha+\beta=\frac{\pi}{2}$
\[
\sin\beta = \cos\alpha,\qquad
\cos\beta = \sin\alpha
\]

Anti-complementair: $|\alpha-\beta|=\frac{\pi}{2}$
\[
\sin\beta = \cos\alpha,\qquad
\cos\beta = -\sin\alpha
\]

Supplementair: $\alpha+\beta=\pi$
\[
\sin(\pi-\alpha)=\sin\alpha,\qquad
\cos(\pi-\alpha)=-\cos\alpha
\]
\[
\tan(\pi-\alpha)=-\tan\alpha,\qquad
\cot(\pi-\alpha)=-\cot\alpha
\]

Anti-supplementair: $|\alpha-\beta|=\pi$
\[
\sin(\pi+\alpha)=-\sin\alpha,\qquad
\cos(\pi+\alpha)=-\cos\alpha
\]
\[
\tan(\pi+\alpha)=\tan\alpha,\qquad
\cot(\pi+\alpha)=\cot\alpha
\]

\section*{4.4 Som-, verschil- en dubbele hoekformules}

Som:
\[
\sin(\alpha+\beta)=\sin\alpha\cos\beta+\cos\alpha\sin\beta
\]
\[
\cos(\alpha+\beta)=\cos\alpha\cos\beta-\sin\alpha\sin\beta
\]
\[
\tan(\alpha+\beta)=\frac{\tan\alpha+\tan\beta}{1-\tan\alpha\tan\beta}
\]

Verschil:
\[
\sin(\alpha-\beta)=\sin\alpha\cos\beta-\cos\alpha\sin\beta
\]
\[
\cos(\alpha-\beta)=\cos\alpha\cos\beta+\sin\alpha\sin\beta
\]
\[
\tan(\alpha-\beta)=\frac{\tan\alpha-\tan\beta}{1+\tan\alpha\tan\beta}
\]

Dubbele hoek:
\[
\sin(2\alpha)=2\sin\alpha\cos\alpha
\]
\[
\cos(2\alpha)=\cos^2\alpha-\sin^2\alpha
\]
\[
\tan(2\alpha)=\frac{2\tan\alpha}{1-\tan^2\alpha}
\]

\section*{4.5 Rechthoekige driehoek}

\[
\sin\alpha=\frac{a}{c},\quad
\cos\alpha=\frac{b}{c},\quad
\tan\alpha=\frac{a}{b}
\]

Pythagoras:
\[
a^2+b^2=c^2
\]

\section*{4.6 Willekeurige driehoek}

Sinusregel:
\[
\frac{a}{\sin\alpha}=\frac{b}{\sin\beta}=\frac{c}{\sin\gamma}
\]

Cosinusregel:
\[
a^2=b^2+c^2-2bc\cos\alpha
\]

Oppervlakte:
\[
A=\frac{1}{2}bc\sin\alpha
\]

\end{document}
